\documentclass[11pt,a4paper]{article}

% --- PACOTES DE MARGENS E FONTE ---
\usepackage[
  a4paper,
  top=3cm,
  bottom=2cm,
  left=3cm,
  right=2cm
]{geometry}

\usepackage{fontspec}
\setmainfont{Arial}

\usepackage{setspace}
\usepackage[brazil]{babel}

\begin{document}

% ==============================================================================
% INÍCIO DA CAPA
% ==============================================================================
\begin{titlepage}
  \begin{center}
    \singlespacing % Capas costumam usar espaçamento simples

    % 1. INSTITUIÇÃO
    \textbf{UNIVERSIDADE DE SÃO PAULO}\\
    \textbf{INSTITUTO DE MATEMÁTICA E ESTATÍSTICA}\\
    \textbf{DEPARTAMENTO DE CIÊNCIA DA COMPUTAÇÃO}\\[4cm]

    % 2. AUTOR
    {\large \textbf{Rafael Vieira de Carvalho}\\[5cm]}

    % 3. TÍTULO E SUBTÍTULO
    {\LARGE \textbf{Divisão de nós em uma max-tree}}\\[0.5cm]
    {\large Atualização dinâmica de max-trees sob transformações
    anti-extensivas locais}\\[4cm]

    {Orientador: Prof. Dr. Ronaldo F. Hashimoto \\
    Coorientador: Prof. Dr. Wonder A. L. Alves}
    
    \vfill % O \vfill empurra o resto para a base

    % 4. LOCAL E ANO
    \textbf{São Paulo}\\
    \textbf{2026}

  \end{center}
\end{titlepage}
% ==============================================================================
% FIM DA CAPA
% ==============================================================================

% Define espaçamento duplo para o restante do documento
\doublespacing

% Inicia a numeração de páginas a partir daqui (página 2 ou reiniciando em 1)
\newpage
\pagenumbering{arabic}

\section{Resumo}

\section{Introdução e Justificativas}
O texto do seu artigo ou projeto de pesquisa começa aqui...

\section{Objetivos}
Colocar aqui perguntas de pesquisa

\section{Metodologia}

\section{Plano de trabalho e cronograma}

\section{}

\section{bibliografia}


\end{document}